\documentclass[10pt,a4paper]{amsart}
\usepackage[margin=19mm]{geometry}
\usepackage{amsmath,amssymb,mathtools}
\usepackage[scr=rsfs,cal=euler]{mathalfa}
\usepackage{xfrac}
\usepackage[unicode,hidelinks,bookmarks=false]{hyperref}
\newcommand{\ie}{i.e.\ }
\newcommand{\cf}{cf.\ }
\newcommand{\resp}{respectively\ }
\newcommand{\Subsection}{\subsection}
\providecommand{\nobreakdash}{\nobreak\hbox{-}}
\setlength{\emergencystretch}{2em}
\multlinegap0pt
\usepackage{bm}
\usepackage{bbm}
\usepackage{dsfont}
\usepackage{xspace}
\usepackage{cancel}
\usepackage{mathtools}
\usepackage{cleveref}
\usepackage{amsthm}
\makeatletter
\@ifundefined{lemma}{\newtheorem{lemma}{Lemma}}{}
\@ifundefined{theorem}{\newtheorem{theorem}{Theorem}}{}
\makeatother
\makeatletter
\newcommand{\R}{\mathbb R}
\newcommand{\X}{\mathcal X}
\newcommand{\Y}{\mathcal Y}
\newcommand{\lsa}{\emph{lsa}\xspace}
\@ifundefined{reflem}{\newenvironment{reflem}[1]{%
  \renewcommand{\reptile}{Lemma \ref{#1}}%
  \begin{repinner}
}{%
  \end{repinner}
}}{}
\@ifundefined{refthm}{\newenvironment{refthm}[1]{%
  \renewcommand{\reptile}{Theorem \ref{#1}}%
  \begin{repinner}
}{%
  \end{repinner}
}}{}
\newcommand{\usa}{\emph{usa}\xspace}
\makeatother
\title[Borel selection of dominating hyperplanes]{Borel selection of dominating hyperplanes}
\author{Eugenio Clerico}
\date{Source: arXiv:2603.19973v2 (CC BY 4.0)}
\begin{document}
\maketitle

\noindent\emph{Source and licence.} Borel selection of dominating hyperplanes. Version arXiv:2603.19973v2 (CC BY 4.0), distributed under the Creative Commons Attribution 4.0 International licence, as recorded in the arXiv licence metadata for this exact version and shown on the abstract page https://arxiv.org/abs/2603.19973v2. Full rights notice: licenses/arxiv-2603.19973-CC-BY-4.0.txt.\par\smallskip

\noindent\textit{Editorial scope.} Counted unit: the Theorem whose source label is \texttt{thm:main} in the article (source0.tex lines 122--125, source label \texttt{thm:main}), delivered together with the article's own complete original proof of it (source0.tex lines 126--156, one \texttt{proof} environment of 31 lines). No mathematics is written, repaired or re-derived: statement and proof are the article's own text line for line. Required context retained uncounted: lemma:sandwich (lines 94--96); lemma:novikov (lines 110--112). Every definition, display and label of those lines is retained unchanged. Omitted from this package and not redistributed: 1--93, 97--109, 113--121, 157--254. Nothing inside a retained excerpt is omitted or altered: every retained body is delivered exactly as the article's lines are written, so no omission receipt of any kind accompanies this package. Citations: the retained excerpts carry no citation command at all, so no bibliography entry of the article is transcribed and no reference is redistributed. Reference adaptation: none was needed and none was made; every reference inside the delivered unit resolves inside it, so no unresolved reference or citation is produced. Printed slips kept literal: no printed word, symbol, hypothesis or number of the counted statement or of its proof was altered. Layout and class: the article's own class is \texttt{article} with packages including inputenc, amsfonts, amsmath, amsthm, amssymb, xcolor, hyperref, cleveref, natbib, xspace, comment; the wrapper is the lane's pinned \texttt{amsart} editorial wrapper at 10pt on A4, which restates the theorem-like environments the retained text needs and adds the article's own macro definitions that the retained text needs (\texttt{\textbackslash R}, \texttt{\textbackslash X}, \texttt{\textbackslash Y}, \texttt{\textbackslash lsa}, \texttt{\textbackslash usa}), with extra wrapper packages (bm, bbm, dsfont, xspace, cancel, mathtools, cleveref). The delivered numbering is the unit's own. Assets: no retained span contains a figure environment, a graphics-inclusion command, a PDF-page inclusion, a table or a TikZ picture; no image file is copied into this package, the package declares no asset dependency and no image file is redistributed with it. Licence: version arXiv:2603.19973v2 of this article is distributed under the Creative Commons Attribution 4.0 International licence, as recorded in the arXiv licence metadata for this exact version and shown on the abstract page https://arxiv.org/abs/2603.19973v2. A full rights notice is delivered at licenses/arxiv-2603.19973-CC-BY-4.0.txt. Characters: every retained excerpt is pure ASCII, so the delivered engine, XeLaTeX with its default Latin Modern text font, reports no missing character.
\begin{lemma}[Analytic sandwich]\label{lemma:sandwich}
Let $\X$ be a standard Borel space, let $u:\X\to\R$ be \usa and $l:\X\to\R$ be \lsa. If $u\leq l$, pointwise on $\X$, then there exists a Borel function $f:\X\to\R$ such that $u\leq f\leq l$, pointwise on $\X$.
\end{lemma}

\begin{lemma}\label{lemma:novikov}
Let $\X$ be a standard Borel space, and let $u:\X\to\R$ be \usa. Then, there exists a Borel map $f:\X\to\R$ such that $f\geq u$ pointwise on $\X$.
\end{lemma}

\begin{theorem}[Hyperplane selection]\label{thm:main}
    Let $\X$ be a standard Borel space and $\Y\subseteq\R^n$ a Borel set. Let $f:\X\times \Y\to\R$ be \usa. Assume that there exist maps $b:\X\to\R^n$ and $c:\X\to\R$ such that $f(x,y)\leq b(x)\cdot y  +c(x)$, for every $x\in\X$ and $y\in \Y$. Then, there exist Borel measurable maps $B:\X\to\R^n$ and $C:\X\to\R$ such that
    $f(x,y)\leq B(x)\cdot y  +C(x)$, for every $x\in\X$ and $y\in\Y$.
\end{theorem}

\begin{proof}
    Without loss of generality, we may assume $\Y = \mathbb R^n$. Indeed, if $\Y \subsetneq \mathbb R^n$, we extend $f$ to $\tilde{f}: \X \times \mathbb R^n \to \R$ by setting $\tilde{f}(x,y) = f(x,y)$ if $y \in \Y$ and $\tilde{f}(x,y) = -\|y\|^2$ (where $\|\cdot\|$ is the Euclidean norm) if $y \notin \Y$. Since $\Y$ is a Borel set, $\tilde{f}$ is \usa. We now verify that $\tilde{f}$ satisfies the affine domination hypothesis. For $y \in \Y$, we have $\tilde{f}(x,y) = f(x,y) \leq b(x)\cdot y + c(x)$ by assumption. On the other hand, for $y \notin \Y$, we use $-\|y\|^2 \leq \|y + b(x)/2\|^2 - \|y\|^2 =  b(x)\cdot y + \|b(x)\|^2/4$. Thus, $\tilde{f}(x,y) \le b(x)\cdot y + \tilde{c}(x)$, for all $x\in\X$ and $y\in\R^n$, where $\tilde{c}(x) = \max(c(x), \|b(x)\|^2/4)$. Therefore, we now proceed assuming $\Y = \mathbb R^n$.

    We prove the result by induction on the dimension. If $n=0$, then $\R^0 = \{0\}$ and the hypothesis  reduces to $f(x,0)\leq c(x)$ for every $x\in\X$. Since $f$ is \usa, the section $x\mapsto f(x,0)$ is also a \usa map from $\X$ to $\R$. By \Cref{lemma:novikov}, there is a real-valued $C:\X\to\R$ such that  $f(x,0)\leq C(x)$ for all $x\in\X$. 

    Now, let $n\geq 1$ and assume that the theorem holds up to $n-1$. We let $f:\X\times \R^n\to\R$  satisfy the affine domination hypothesis with $b$ and $c$.  
    We partition $\R^n$ into $\R^n_> = \{y\in \R^n\,:\,y_n>0\}$, $\R^n_<=\{y\in \R^n\,:\,y_n<0\}$, and $\R^n_0 = \{y\in \R^n\,:\,y_n=0\}$. Define the set 
    $$\Gamma = \left\{(\tilde y, y, y')\in \R^n_0\times \R^n_>\times \R^n_<\,:\,\frac{y_{n} y' - y'_{n} y}{y_{n}-y'_{n}} = \tilde y\right\}\,.$$ This is a Borel set, and so the function $g:\X\times \R^n_0\times \R^n_>\times \R^n_<\to[-\infty,\infty)$, defined as 
    $$g(x,\tilde y, y, y') = \begin{cases}
    \frac{y_{n}f(x,y') - y'_{n}f(x,y)}{y_{n}-y'_{n}} & \text{if $(\tilde y, y, y')\in\Gamma$};\\-\infty & \text{otherwise,}\end{cases}$$
    is \usa (note that $y_{n}>0$ and $y'_{n}<0$). For any $(x, \tilde y, y, y')\in \X\times \R^n_0\times \R^n_>\times \R^n_<$,$$g(x, \tilde y, y, y') \leq \frac{y_{n}  \big(b(x)\cdot y'+c(x)\big) - y'_{n} \big(b(x)\cdot y+c(x)\big)}{y_{n}-y'_{n}} =   b(x) \cdot\tilde y + c(x)\,,$$  as the left side is finite only if $(\tilde y, y,y')\in\Gamma$. Since the upper-bound on $g$ is always finite, we can define $h:\X\times \R^n_0\to[-\infty,\infty)$ as  
    $$h(x,\tilde y) = \sup_{y\in \R^n_>,y'\in \R^n_<} g(x, \tilde y, y, y')\,,$$ 
    which is \usa, as the supremum of a \usa function over an analytic set, and satisfies $h(x,\tilde y)\leq b(x)\cdot \tilde y + c(x)$, for all $(x,\tilde y)\in\X\times \R^n_0$.

    We can now define $F:\X\times \R^n_0\to \R$ as
    $$F(x, \tilde y) = \max\big(h(x,\tilde y),f(x,\tilde y)\big)\,.$$ $F$ is the maximum between two \usa functions, hence it is \usa. Moreover, since both $h$ and $f$ are dominated by $(x,\tilde y)\mapsto b(x)\cdot\tilde y + c(x)$, we have that
    $$F(x,\tilde y) \leq b(x)\cdot \tilde y + c(x) = \sum_{i=1}^{n-1} b_i(x)\tilde y_i + c(x)\,,$$ for every $\tilde y\in \R^n_0$ and $x\in\X$, where we used that $\tilde y_{n} = 0$ as $\tilde y\in \R^n_0$. By the inductive hypothesis, since  $\R^n_0$ can be canonically identified with  $\R^{n-1}$, we have that there exist Borel functions (from $\X$ to $\R$) $B_1,\dots, B_{n-1}$, and $C$, such that
    $$F(x, \tilde y) \leq \sum_{i=1}^{n-1} B_i(x)\tilde y_i + C(x)\,,$$
    for every $x\in\X$ and $\tilde y\in \R^n_0$. 

    Fix $y\in \R^n_>$ and $y'\in \R^n_<$, and let $\tilde y = \frac{y_{n}y' - y_{n}'y}{y_{n}-y'_{n}}$. Clearly, $\tilde y\in \R^n_0$ and $(\tilde y, y,y')\in\Gamma$. In particular,
    $$\frac{y_{n}f(x,y') - y'_{n}f(x,y)}{y_{n}-y'_{n}}= g(x, \tilde y, y, y')\leq F(x,\tilde y) \leq \sum_{i=1}^{n-1}B_i(x)\frac{y_{n}y_i' - y_{n}'y_i}{y_{n}-y'_{n}} + C(x)\,.$$
    Rearranging (and using that $y_{n}-y'_{n}>0$), we get
    $$\frac{f(x,y)-\sum_{i=1}^{n-1} B_i(x)y_i -C(x)}{y_{n}}\leq \frac{f(x,y')-\sum_{i=1}^{n-1} B_i(x)y'_i-C(x)}{y'_{n}}\,.$$
    This holds for every $x\in\X$, $y\in \R^n_>$, and $y'\in \R^n_<$. We now define the functions $U$ and $L$, from $\X$ to $[-\infty,\infty]$, as
    \begin{align*}U(x) &= \sup_{y\in \R^n_>}\frac{f(x,y)-\sum_{i=1}^{n-1} B_i(x)y_i -C(x)}{y_{n}}\,;\\ L(x) &= \inf_{y'\in \R^n_<}\frac{f(x,y')-\sum_{i=1}^{n-1} B_i(x)y'_i-C(x)}{y'_{n}}\,.\end{align*} Since $\R^n_<$ and $\R^n_>$ are non-empty, $L(x)<\infty$ and $U(x)>-\infty$ for any $x\in\X$. As we have $U\leq L$ pointwise, we can view $U$ and $L$ as functions from $\X$ taking values in $\R$. Note that $U$ is \usa and $L$ is \lsa, since they are respectively the supremum and the infimum on analytic domains of \usa and \lsa functions. In particular, the analytic sandwich lemma (\Cref{lemma:sandwich}) applies and we can find a Borel map $B_n:\X\to\R$ such that $U(x)\leq B_n(x)\leq L(x)$, for all $x\in\X$. 
    
    We now define $B:\X\to\R^n$ as the Borel map with components $B_1,\dots, B_n$. We have  obtained that, for any $x\in\X$ and $y\in \R^n_>\cup \R^n_<$, 
    $$f(x,y) \leq B(x)\cdot y + C(x)\,.$$ We are only left to check the case $y\in \R^n_0$. However, in such case we simply have that
    $$f(x,y)\leq F(x,y)\leq \sum_{i=1}^{n-1} B_i(x) y_i +C(x) = B(x)\cdot y + C(x)\,,$$ since $y_{n}=0$.
\end{proof}

\end{document}
